% !TeX root = ../../Book2.tex
% 纲要标题：### 19.2 Goldie 理论
% 纲要位置：计划纲要.md，第 1676 行。

\section{Goldie 理论}
\label{b2:sec:1902}

Ore 条件逐个要求方程 \(at=sb\) 可解，直接核验往往困难。Goldie 理论用环的右理想结构保证这些方程有解，并进一步说明所得分式环为何半单。本节明确改用右模：一致维数、子模和零化子链均按所写的作用方向理解。证明先建立本质右理想中正则元素的存在性，再进行局部化。

% 纲要要点（供目录核对）：
%
% * 一致维数及零化子条件；以整数、向量空间和截断多项式模区别本质、一致、长度，解释有限本质一致直和的构造与逆向证明
% * 半素右 Goldie 环的分式环；分别追踪半素性、右零化子升链条件及有限一致维数，补清本质理想含正则元、Ore 乘子的原像和右半单转为左半单的步骤
% * 定理假设与交换整环的比较；正文证明整环右 Ore 等价于右正则模一致，正式证明上三角环中的本质理想无正则元，区分矩阵的行模数与基域维数
%

\subsection{一致维数与零化子条件}
\label{b2:subsec:1902:01}

\begin{definition}\label{b2:def:essential-uniform-dimension}
设 \(R\) 为含幺环，\(M\) 为幺性右 \(R\) 模，\(E\subseteq M\) 为子模。若每个非零子模 \(N\subseteq M\) 都满足 \(N\cap E\ne0\)，则称 \(E\) 为 \(M\) 的\term[benzhi-zimo]{本质子模}[essential submodule]。若非零幺性右 \(R\) 模 \(U\) 的任意两个非零子模交非零，则称 \(U\) 为\term[yizhi-mo]{一致模}[uniform module]；等价地，\(U\) 的每个非零子模都在 \(U\) 中本质。

若 \(M\) 的非零子模 \(N_1,\ldots,N_r\) 的和为内部直和，则称它们独立。定义 \(M\) 的\term[yizhi-weishu]{一致维数}[uniform dimension] \(\operatorname{udim}M\) 为这类有限独立族大小的上确界，允许为无穷；零模的值为零。一致维数有限时也称为 \term[Goldie-zhi]{Goldie 秩}[Goldie rank]。
\end{definition}
\symidx{\(\operatorname{udim}M\)：右模的一致维数}

本质包含要求每个非零子模都与给定子模相交，并不要求给定子模等于全模。例如 \(2\mathbb Z\) 在 \(\mathbb Z\) 中本质，因为任意非零子模 \(d\mathbb Z\) 都与它有非零公共元素 \(2d\)。但域 \(k\) 上向量空间 \(V\) 的真子空间 \(W\) 不本质：取 \(v\notin W\)，则 \(kv\cap W=0\)。

\ProseParagraphBreak
一致模不能同时容纳两个交零的非零子模，所以对非零模 \(U\)，一致性等价于 \(\operatorname{udim}U=1\)。对除环上的右向量空间，这恰好是维数为一的条件：一维空间只有零子空间与自身，而两个线性无关向量张成的直线交于零。简单模因此一致，但一致模未必简单。例如交换整环 \(R\) 的右正则模一致：两个非零理想分别取非零元素 \(a,b\)，非零乘积 \(ab\) 落在两者的交中。整数模因而一致，虽然有无限严格降链 \(\mathbb Z\supsetneq2\mathbb Z\supsetneq4\mathbb Z\supsetneq\cdots\)，并不具有有限长度。

\ProseParagraphBreak
一致维数数的是同时独立的直和项，合成长度则还计算这些项内部的简单层。取域 \(k\)、整数 \(m\ge1\) 及 \(R=k[t]/(t^m)\)，把 \(t\) 的剩余类仍记为 \(t\)。右理想与 \(k[t]\) 中包含 \((t^m)\) 的理想对应；\(k[t]\) 为主理想整环，这些理想的生成元整除 \(t^m\)，故可取为 \(t^i\)。所以 \(R\) 的右理想恰为 \(t^iR\)、\(0\le i\le m\)。它们彼此可比，任意两个非零右理想的交非零，因此 \(\operatorname{udim}R_R=1\)。另一方面，链
\[
0=t^mR\subsetneq t^{m-1}R\subsetneq\cdots\subsetneq tR\subsetneq R
\]
的每个相邻商 \(t^iR/t^{i+1}R\) 都由 \(t^i\) 的类张成，且 \(t\) 在该商上作用为零，所以同构于简单右模 \(R/tR\cong k\)。这给出长度为 \(m\) 的合成列，故 \(\ell(R_R)=m\)。

\begin{lemma}\label{b2:lem:essential-tools}
设 \(R\) 为含幺环，所涉右 \(R\) 模均为幺性模。若 \(M\) 的子模 \(E\subseteq N\subseteq M\) 满足 \(E\) 在 \(N\) 中本质、\(N\) 在 \(M\) 中本质，则 \(E\) 在 \(M\) 中本质。若 \(E\subseteq N\) 是右 \(R\) 模中的本质包含，\(f:M\rightarrow N\) 是右 \(R\) 模同态，则 \(f^{-1}(E)\) 在 \(M\) 中本质。对有限族右 \(R\) 模中的本质包含 \(E_i\subseteq M_i\)，还有本质包含 \(\bigoplus_iE_i\subseteq\bigoplus_iM_i\)。
\end{lemma}
\begin{proof}
若 \(E\) 本质于 \(N\)，而 \(N\) 本质于 \(M\)，任意非零 \(L\subseteq M\) 先与 \(N\) 非零相交，再与 \(E\) 非零相交，故传递性成立。

对非零子模 \(L\subseteq M\)，若 \(f(L)=0\)，则 \(L\subseteq f^{-1}(E)\)；否则 \(f(L)\cap E\ne0\)，取一个非零像的原像 \(x\in L\)，便有 \(0\ne x\in L\cap f^{-1}(E)\)。这证明原像的本质性。

在有限直和中，第 \(i\) 个投影的原像 \(\pi_i^{-1}(E_i)\) 由上段本质。有限个本质子模的交仍本质：先与第一个相交，再逐个与其余相交，每一步保持非零。它们的交恰为 \(\bigoplus_iE_i\)，得到最后结论。
\end{proof}

\begin{lemma}\label{b2:lem:uniform-submodule-existence}
设 \(R\) 为含幺环，\(M\) 为非零幺性右 \(R\) 模。若 \(M\) 的一致维数有限，或 \(M\) 为 Noether 模，则 \(M\) 的每个非零子模都含有一致子模。
\end{lemma}
\begin{proof}
假设某个非零子模 \(N\subseteq M\) 不含一致子模。令 \(N_0=N\)。由于 \(N_0\) 不一致，可取非零子模 \(N_1,W_0\subseteq N_0\)，使 \(N_1\cap W_0=0\)。\(N_1\) 也不含一致子模，故可继续在 \(N_1\) 中取非零子模 \(N_2,W_1\)，使 \(N_2\cap W_1=0\)，依此类推。每一步都有 \(N_{i+1}\oplus W_i\subseteq N_i\)，归纳可知 \(W_0,W_1,\ldots\) 独立。

若 \(\operatorname{udim}M<\infty\)，任意多个 \(W_i\) 构成的有限独立族便与一致维数的上界矛盾。若 \(M\) 为 Noether 模，则 \(W_0\subsetneq W_0\oplus W_1\subsetneq\cdots\) 是子模的无限严格升链，同样矛盾。
\end{proof}

\begin{proposition}\label{b2:prop:uniform-dimension-structure}
设 \(R\) 为含幺环，\(M\) 为非零幺性右 \(R\) 模。若 \(\operatorname{udim}M=d<\infty\)，则 \(M\) 含有本质子模 \(E=U_1\oplus\cdots\oplus U_d\)，其中各 \(U_i\) 是一致右 \(R\) 模。反过来，只要 \(M\) 含有这样的有限本质直和，就有 \(\operatorname{udim}M=d\)。若子模 \(N\subseteq M\) 且 \(M\) 一致维数有限，则
\[
\operatorname{udim}N=\operatorname{udim}M
\quad\Longleftrightarrow\quad N\text{ 在 }M\text{ 中本质}.
\]
每个 Noether 右 \(R\) 模都有有限一致维数。
\end{proposition}
\begin{proof}
先设 \(\operatorname{udim}M=d<\infty\)。由\cref{b2:lem:uniform-submodule-existence}取一致子模 \(U_1\)。若已取的有限直和 \(U_1\oplus\cdots\oplus U_r\) 不本质，就有非零子模与它交于零，在其中再取一致子模 \(U_{r+1}\)。一致维数限制新增项数，所以过程在有限步后停止，得到本质直和 \(E\)。

现在证明：本质直和 \(E=\bigoplus_{i=1}^dU_i\) 存在时，任意独立非零族 \(N_1,\ldots,N_s\subseteq M\) 都满足 \(s\le d\)。先以 \(N_j\cap E\) 代替 \(N_j\)，这些交仍非零且独立，故可假设它们在 \(E\) 中。对 \(d\) 归纳，\(d=0\) 时没有非零子模。若 \(s>0\)，令 \(W=N_2\oplus\cdots\oplus N_s\)。它与非零 \(N_1\) 交于零，所以不在 \(E\) 中本质。要减少一致直和的项数，可以找一个与 \(W\) 交零的 \(U_i\)，再将它商去；关键是证明这样的 \(U_i\) 必存在。如果每个 \(U_i\cap W\) 都非零，一致性及\cref{b2:lem:essential-tools}将使这些交的直和在 \(E\) 中本质，而该直和包含于 \(W\)，迫使 \(W\) 与 \(N_1\) 非零相交，矛盾。因此某个 \(U_i\cap W=0\)。商映射 \(E\rightarrow E/U_i\) 在 \(W\) 上单射，故其余 \(s-1\) 项在 \(E/U_i\cong\bigoplus_{j\ne i}U_j\) 中仍独立。归纳给出 \(s-1\le d-1\)。原来的 \(U_i\) 又给出大小 \(d\) 的独立族，故一致维数恰为 \(d\)。因此，前段构造的本质直和恰有 \(d\) 项。

若 \(M\) 为 Noether 模，\cref{b2:lem:uniform-submodule-existence}同样保证每次都能取出一致子模。若抽取过程不停，所得部分直和就构成严格升链 \(U_1\subsetneq U_1\oplus U_2\subsetneq\cdots\)，违反第十二章的子模升链条件。因此过程在有限步后停止，所得本质直和有有限项数；由刚证的逆向结论，\(M\) 的一致维数有限。

最后，若 \(N\) 本质，把 \(M\) 中任意最大独立族的每项与 \(N\) 相交，得到 \(N\) 中同样大的独立族，所以两维数相等。若 \(N\) 不本质，选 \(0\ne W\subseteq M\) 与 \(N\) 交零。将它加入 \(N\) 的一个最大独立族，得到 \(\operatorname{udim}M\ge\operatorname{udim}N+1\)，故两者不等。
\end{proof}

对 \(X\subseteq R\)，分别记
\[
\begin{aligned}
\operatorname{ann}_r(X)&=\{\,a\in R:xa=0\text{ 对所有 }x\in X\,\},\\
\operatorname{ann}_{\ell}(X)&=\{\,a\in R:ax=0\text{ 对所有 }x\in X\,\}.
\end{aligned}
\]
它们分别为右理想与左理想，称为右零化子与左零化子。
\symidx{\(\operatorname{ann}_r(X)\)：子集的右零化子}
\symidx{\(\operatorname{ann}_{\ell}(X)\)：子集的左零化子}

\begin{definition}\label{b2:def:semiprime-right-goldie}
设 \(R\) 为含幺环。若 \(R\) 没有非零的平方为零的双边理想，则称 \(R\) 为\term[bansu-huan]{半素环}[semiprime ring]。对 \(X\subseteq R\)，记 \(\operatorname{ann}_r(X)=\{\,r\in R\mid xr=0\text{ 对所有 }x\in X\,\}\)。若右正则模 \(R_R\) 一致维数有限，且这些右零化子组成的每条升链都最终稳定，则称 \(R\) 为\term[you-Goldie-huan]{右 Goldie 环}[right Goldie ring]。
\end{definition}

半素也等价于没有非零幂零双边理想：若 \(I^m=0\) 且 \(m\ge2\) 最小，则 \(I^{m-1}\ne0\) 而它的平方为零。这里排除的是双边理想，不是幂零元素；例如对域 \(k\)，单环 \(M_2(k)\) 半素，却有 \(E_{12}\ne0\)、\(E_{12}^2=0\)。半素环中的非零左理想或右理想，其平方也不能为零。比如左理想 \(L\) 若满足 \(L^2=0\)，双边理想 \(LR\) 满足
\((LR)^2\subseteq L^2R=0\)，故 \(LR=0\)，进而 \(L=0\)；右理想采用 \(RL\) 的对应论证。这里零化子升链条件只针对这些特殊右理想，比全部右理想都满足升链条件弱。

\subsection{半素右 Goldie 环的分式环}
\label{b2:subsec:1902:02}

下面要从半素右 Goldie 条件推出每个本质右理想都包含正则元素。半素性使非零一侧理想中能找到乘积非零的元素，右零化子的升链条件允许选择不能继续扩大的零化子；有限一致维数则限制可以加入的独立子模数目，使寻找单射左乘作用的过程在有限步后完成。三个条件将分别用于这些具体论证。

\begin{lemma}\label{b2:lem:essential-annihilator-zero}
设 \(R\) 为非零含幺半素右 Goldie 环，\(E\subseteq R_R\) 为本质右理想，则其左零化子 \(\operatorname{ann}_{\ell}(E)=\{\,r\in R\mid re=0\text{ 对所有 }e\in E\,\}\) 为零。
\end{lemma}
\begin{proof}
若存在本质右理想 \(E\) 满足 \(\operatorname{ann}_{\ell}(E)\ne0\)，取非零左理想 \(L=\operatorname{ann}_{\ell}(E)\)，则 \(\operatorname{ann}_r(L)\supseteq E\)，因而也是本质右理想。故下列右零化子的集合非空：
\[
\mathcal A=
\bigl\{\operatorname{ann}_r(L):
0\ne L\subseteq R\text{ 为左理想，且 }\operatorname{ann}_r(L)\text{ 在 }R_R\text{ 中本质}\bigr\}.
\]
沿用\cref{b2:prop:chain-maximal-minimal}证明中的升链论证，若 \(\mathcal A\) 没有极大元，就可在其中逐次严格增大，得到一条永不稳定的右零化子升链，违反右 Goldie 条件。因此可取 \(A=\operatorname{ann}_r(L)\in\mathcal A\) 为按包含关系极大的元素。

半素性给出 \(L^2\ne0\)，取 \(x,y\in L\) 使 \(xy\ne0\)。\(yR\) 是非零右理想，本质性使其中有非零元素 \(ya\in A\)。因为 \(x\in L\)、\(ya\in\operatorname{ann}_r(L)\)，有 \(xya=0\)，而 \(ya\ne0\)；也就是说，右乘子 \(a\) 零化了 \(xy\)，却没有零化 \(y\)。所以
\[
A\subseteq\operatorname{ann}_r(y)\subsetneq\operatorname{ann}_r(xy).
\]
这里 \(y,xy\in L\)，所以两者都包含 \(A\)；并且
\[
\operatorname{ann}_r(y)=\operatorname{ann}_r(Ry),\qquad
\operatorname{ann}_r(xy)=\operatorname{ann}_r(Rxy).
\]
由于 \(Ry,Rxy\) 均为非零左理想，而包含 \(A\) 的右理想仍本质，这两个零化子都属于 \(\mathcal A\)。\(A\) 的极大性要求它们都等于 \(A\)，与严格包含矛盾。
\end{proof}

\begin{lemma}\label{b2:lem:one-sided-regular-goldie}
设 \(R\) 为非零含幺半素右 Goldie 环，\(s\in R\)。若 \(\operatorname{ann}_r(s)=\{\,r\in R\mid sr=0\,\}=0\)，则 \(s\) 双侧正则，且 \(sR\) 为 \(R_R\) 的本质右理想。
\end{lemma}
\begin{proof}
左乘 \(s\) 给出右模同构 \(R_R\cong sR\)，所以两者一致维数相同。由\cref{b2:prop:uniform-dimension-structure}，\(sR\) 在 \(R\) 中本质。上一引理给出 \(\operatorname{ann}_{\ell}(sR)=0\)，而它正是 \(\operatorname{ann}_{\ell}(s)\)。两侧零化子均为零，所以 \(s\) 正则。
\end{proof}

\begin{lemma}\label{b2:lem:goldie-ideal-regular-element}
设 \(R\) 为非零含幺半素右 Goldie 环，则每个右理想 \(N\subseteq R\) 都含有某个 \(v\in N\)，使 \(\operatorname{ann}_r(v)\cap N=0\)。因此每个本质右理想都含有双侧正则元素。
\end{lemma}
\begin{proof}
先设 \(U\) 为非零一致右理想。由半素性，取 \(x,y\in U\) 使 \(xy\ne0\)。若 \(W=\operatorname{ann}_r(x)\cap U\ne0\)，则 \(W\) 在 \(U\) 中本质。对所有满足 \(ya\in W\) 的右乘子 \(a\)，都有 \(xya=0\)；为了把这些乘子组成的右理想与本质性联系起来，考虑右线性映射 \(\lambda_y:R\rightarrow U\)、\(a\mapsto ya\)。这些乘子恰好组成 \(\lambda_y^{-1}(W)\)，由\cref{b2:lem:essential-tools}它是本质右理想，却被非零元素 \(xy\) 从左零化，违反\cref{b2:lem:essential-annihilator-zero}。故 \(W=0\)，一致右理想中的结论成立。

若 \(N\ne0\)，它含有一致右理想，上一段因而保证存在满足
\[
v\in V\subseteq N,\qquad \operatorname{ann}_r(v)\cap V=0
\]
的非零子模 \(V\) 及其中的元素 \(v\)。选择一对 \((V,v)\)，使 \(\operatorname{udim}V\) 最大；这些维数是被 \(\operatorname{udim}R_R\) 限制的正整数，所以最大值存在。假设仍有 \(\operatorname{ann}_r(v)\cap N\ne0\)，在其中取一致右理想 \(U\)，再由上段取 \(u\in U\) 使 \(\operatorname{ann}_r(u)\cap U=0\)。因为 \(U\subseteq\operatorname{ann}_r(v)\)，有 \(U\cap V=0\)。

令 \(x=u+v\)，是为了在保留 \(v\) 对 \(V\) 的单射作用时，用 \(u\) 补上它在 \(U\) 上的零作用。这里 \(U,V\) 都是右理想，所以 \(u(U\oplus V)\subseteq U\)、\(vV\subseteq V\)，又有 \(vU=0\)。因此左乘 \(x\) 确实给出右模自同态 \(\lambda_x:U\oplus V\rightarrow U\oplus V\)，相对于这个分解写为
\[
\lambda_x=
\begin{pmatrix}
\lambda_u|_U&\lambda_u|_V\\
0&\lambda_v|_V
\end{pmatrix}.
\]
两条对角映射按 \(u,v\) 的选择都是单射，左下角则为零，所以相加不会产生新的核。具体地，若 \(z=a+b\)、\(a\in U,b\in V\)，且 \(xz=0\)，先看 \(V\) 分量得到 \(vb=0\)，从而 \(b=0\)；再看 \(U\) 分量得到 \(ua=0\)，从而 \(a=0\)。故 \(\operatorname{ann}_r(x)\cap(U\oplus V)=0\)。但 \(x\in U\oplus V\subseteq N\)，而把 \(V\) 中大小为 \(\operatorname{udim}V\) 的独立族与非零项 \(U\) 并在一起，就使一致维数至少增加一，与选择的最大值矛盾。因此 \(\operatorname{ann}_r(v)\cap N=0\)。\(N=0\) 时取 \(v=0\) 即可。

若 \(N\) 本质，这个零交迫使右理想 \(\operatorname{ann}_r(v)\) 为零。由\cref{b2:lem:one-sided-regular-goldie}，\(v\) 双侧正则。
\end{proof}

\begin{theorem}[Goldie 定理的充分方向]\label{b2:thm:goldie-sufficiency}
设 \(R\) 为非零含幺半素右 Goldie 环，\(S\subseteq R\) 为全部双侧正则元素，则 \(S\) 满足右 Ore 条件，经典右分式环 \(Q=RS^{-1}\) 为半单 Artin 环。
\end{theorem}
\begin{proof}
一属于 \(S\)，正则元素乘法封闭。给定 \(a\in R,s\in S\)，要解 \(at=sb\)，就是要找到正则的右乘子 \(t\)，使 \(at\in sR\)。这些候选乘子组成右模同态 \(\lambda_a:R_R\rightarrow R_R\)、\(x\mapsto ax\) 的原像 \(\lambda_a^{-1}(sR)\)。由\cref{b2:lem:one-sided-regular-goldie}，\(sR\) 本质，所以\cref{b2:lem:essential-tools}使这个原像也本质。再由\cref{b2:lem:goldie-ideal-regular-element}取其中的正则元素 \(t\)，便有 \(at=sb\) 对某个 \(b\in R\) 成立，得到右 Ore 条件。\cref{b2:thm:right-ore-localization}给出 \(Q\)。

局部化把正则元素变成单位，这将排除 \(Q\) 中的真本质右理想。设 \(E\) 为 \(Q\) 的本质右理想，\(I\) 为 \(R\) 的任意非零右理想，则 \(IQ\) 是非零右 \(Q\) 理想，所以 \(E\cap IQ\) 含有非零元素 \(q\)。由\cref{b2:prop:ore-clearing-denominators}，可取 \(s\in S\) 使 \(qs\in I\)，且 \(s\) 在 \(Q\) 中可逆保证 \(qs\ne0\)。因为 \(E\) 是右理想，\(qs\) 也在 \(E\)，因此 \(E\cap R\) 是 \(R\) 的本质右理想。它含有正则元素 \(t\)，而 \(t^{-1}\in Q\)；右理想 \(E\) 因而含有 \(tt^{-1}=1\)，只能等于 \(Q\)。

要将本质理想的结论变成直和分解，还需为任意右理想找到与它交零、且二者之和本质的子模。任意独立的非零右 \(Q\) 理想与 \(R\) 的交均非零，且作为右 \(R\) 子模仍独立，所以 \(\operatorname{udim}Q_Q\le\operatorname{udim}R_R<\infty\)。对给定右 \(Q\) 理想 \(J\)，在与 \(J\) 交零的右理想中选一致维数最大的 \(K\)：零理想保证候选非空，而这些维数是有上界的非负整数，保证最大值存在。若 \(J\oplus K\) 不本质，取非零右理想 \(L\) 与之交零；则 \(K\oplus L\) 仍与 \(J\) 交零，且将 \(L\) 加入 \(K\) 的最大独立族使一致维数至少增加一，违反所选最大值。故 \(J\oplus K\) 本质。刚证 \(Q\) 没有真本质右理想，所以 \(J\oplus K=Q\)。

因此右正则模 \(Q_Q\) 的每个子模都有补模。此前的\cref{b2:thm:semisimple-module-characterizations}以左模陈述，而半单环也按左正则模定义；把右 \(Q\) 模视为左 \(Q^{\mathrm{op}}\) 模，便可对同一个空间及其子模应用这一刻画，得到
\[
Q_Q\text{ 半单}
\quad\Longleftrightarrow\quad
{}_{Q^{\mathrm{op}}}Q\text{ 半单}
\quad\Longleftrightarrow\quad
Q^{\mathrm{op}}\text{ 为半单环}.
\]
对 \(Q^{\mathrm{op}}\) 应用\cref{b2:thm:artin-wedderburn}，存在有限多个除环 \(D_i\) 和正整数 \(n_i\)，使
\[
Q^{\mathrm{op}}\cong\prod_i M_{n_i}(D_i).
\]
取反环，并以矩阵转置识别 \(M_{n_i}(D_i)^{\mathrm{op}}\cong M_{n_i}(D_i^{\mathrm{op}})\)，便得
\[
Q\cong\prod_i M_{n_i}(D_i^{\mathrm{op}}).
\]
故 \(Q\) 是半单 Artin 环。
\end{proof}

这给出本章所用的\term[Goldie-dingli]{Goldie 定理}[Goldie's theorem]：由半素右 Goldie 条件得到半单经典分式环。本章证明这一充分方向，反向刻画留待进一步阅读。

\begin{corollary}\label{b2:cor:goldie-rank}
设 \(R\) 为非零含幺半素右 Goldie 环，\(S\subseteq R\) 为全部双侧正则元素，\(Q=RS^{-1}\) 为经典右分式环，则 \(\operatorname{udim}R_R=\ell(Q_Q)\)。
\end{corollary}
\begin{proof}
若 \(I_1,\ldots,I_d\) 为 \(R\) 中独立非零右理想，则 \(I_iQ\) 仍独立：一条和为零的关系中的各项，通到共同右分母后成为各 \(I_i\) 中和为零的分子，原独立性使每个分子为零。反过来，\(Q\) 中独立非零右理想与 \(R\) 的交都非零，因为单个非零分式可清分母；这些交仍独立。因此 \(R_R,Q_Q\) 一致维数相同。\(Q_Q\) 是有限个简单右模的直和，简单模一致，\cref{b2:prop:uniform-dimension-structure}说明其一致维数正好等于简单分量个数，也就是长度。
\end{proof}

\subsection{Goldie 假设与交换整环}
\label{b2:subsec:1902:03}

若 \(R\) 右 Noether，则全部右理想满足升链条件，特别是右零化子链稳定；\cref{b2:prop:uniform-dimension-structure}又保证右正则模一致维数有限。因此半素右 Noether 环满足 Goldie 定理。对允许非交换的整环，非零双边理想中的 \(a\ne0\) 满足 \(a^2\ne0\)，所以半素性自动成立，右 Noether 整环因而有经典右分式环；分母包含全部非零元素，使它成为除环，具体见\cref{b2:cor:domain-uniform-ore}。

\ProseParagraphBreak
交换整环甚至不必 Noether：其右正则模一致，而含非零元素的子集零化子为零，故零化子链自动稳定。Goldie 定理便恢复通常的分式域。

\ProseParagraphBreak
在 \(R=M_n(C)\)、\(C\) 为交换整环、\(n\ge1\) 的例子中，正则矩阵恰为行列式非零的矩阵。行列式非零时，伴随矩阵公式及 \(C\) 中的消去律排除两侧零化元；行列式为零时，在分式域中取非零核列，清分母得非零的 \(C\) 系数核列，将它放入一个矩阵的第一列，其余列为零，就得到非零右零化元。因此所有正则矩阵都在 \(Q=M_n(\operatorname{Frac}C)\) 中可逆，而该环的每个元素已可用非零标量矩阵作一个右分母，故它正是经典分式环。由\cref{b2:thm:artin-wedderburn}证明中的行分解，\(Q_Q=E_{11}Q\oplus\cdots\oplus E_{nn}Q\)，每个行理想都是简单右模，故 \(\ell(Q_Q)=n\)。清分母保持独立理想族的同一论证给出 \(\operatorname{udim}R_R=n\)。这里数的是 \(n\) 个简单行模，不是作为底域向量空间的 \(n^2\) 个矩阵坐标。

\ProseParagraphBreak
半素性不能删去。\cref{b2:prop:essential-without-semiprime}将对二阶上三角矩阵环给出一个本质右理想，其中没有正则元素；该环仍满足右 Goldie 条件，其经典分式环也仍是原来的非半单环。

\ProseParagraphBreak
有限一致维数也不能由另外两个条件替代。取域 \(k\)，自由代数 \(R=k\langle x,y\rangle\) 是整环，故半素；其含有非零元素的子集右零化子均为零，所以零化子升链条件成立。但右理想
\[
yR,\ xyR,\ x^2yR,\ldots
\]
独立，因为各自的基字以前缀 \(x^iy\) 开始，这些前缀彼此不会延长成另一个。故一致维数无穷，不能应用 Goldie 定理。这与上一节发现它不满足右 Ore 条件相吻合。

\begin{corollary}\label{b2:cor:domain-uniform-ore}
设 \(R\) 为非零含幺整环，不要求交换，并令 \(S=R\setminus\{0\}\)。则 \(S\) 满足右 Ore 条件，当且仅当幺性右正则模 \(R_R\) 一致。条件成立时，经典右分式环 \(RS^{-1}\) 为除环。
\end{corollary}
\begin{proof}
整环中每个非零元素都双侧正则，\(S\) 含一且乘法封闭。若 \(S\) 满足右 Ore 条件，在任意两个非零右理想 \(I,J\) 中分别取 \(a,s\ne0\)，再取 \(t\in S,b\in R\) 使 \(at=sb\)。整环性保证 \(at\ne0\)，所以 \(0\ne at\in I\cap J\)，右正则模一致。

反过来，若 \(R_R\) 一致，对 \(a,s\ne0\)，两个非零右理想 \(aR,sR\) 的交中有非零元素 \(at=sb\)。这一元素非零，保证 \(t\ne0\)，即 \(t\in S\)。若 \(a=0\)，取 \(t=1,b=0\) 即可。因此对全部 \(a\in R,s\in S\) 都满足右 Ore 条件。

由\cref{b2:thm:right-ore-localization}得到单射 \(R\hookrightarrow Q=RS^{-1}\)。每个非零分式 \(as^{-1}\) 的分子也非零，故 \(a\) 与 \(s\) 都在 \(Q\) 中可逆，其逆为 \(sa^{-1}\)。所以 \(Q\) 是除环。
\end{proof}

\begin{proposition}\label{b2:prop:essential-without-semiprime}
设 \(k\) 为域，\(R=T_2(k)\) 为含幺二阶上三角矩阵环，并令
\[
E=\left\{\begin{pmatrix}0&b\\0&c\end{pmatrix}:b,c\in k\right\}.
\]
则 \(E\) 是 \(R_R\) 的本质右理想，却不含双侧正则元素。环 \(R\) 满足右 Goldie 条件，但不半素；它的正则元素恰为单位，经典右分式环仍为 \(R\)，因而不半单。
\end{proposition}
\begin{proof}
右乘上三角矩阵保持第一列为零，所以 \(E\) 是右理想。给定非零右理想 \(I\)，取
\[
0\ne A=\begin{pmatrix}a&b\\0&c\end{pmatrix}\in I.
\]
若 \(a=0\)，则 \(A\in I\cap E\) 已非零；若 \(a\ne0\)，则 \(AE_{12}=aE_{12}\ne0\) 也属于 \(I\cap E\)。因此 \(E\) 本质。但对每个 \(A\in E\) 都有 \(AE_{11}=0\)，而 \(E_{11}\ne0\)，所以 \(E\) 没有正则元素。

环 \(R\) 作为 \(k\) 空间只有三维，故右理想升链必稳定，即 \(R\) 右 Noether；这保证右零化子升链稳定，并由\cref{b2:prop:uniform-dimension-structure}保证一致维数有限，所以 \(R\) 右 Goldie。严格上三角双边理想 \(kE_{12}\) 非零且平方为零，故 \(R\) 不半素。半单环的 Artin--Wedderburn 分解不允许非零平方为零的双边理想，因此 \(R\) 也不半单。

最后，对上面的 \(A\)，若 \(a,c\ne0\)，就有上三角逆，所以 \(A\) 正则。若 \(a=0\)，则 \(AE_{11}=0\)；若 \(c=0\)，则 \(E_{22}A=0\)。因此正则元素恰为单位。对任意 \(r\in R,s\in R^\times\)，取 \(t=1\)、\(b=s^{-1}r\)，就有 \(rt=sb\)，即满足右 Ore 条件。将已有单位求逆不增加元素，故经典右分式环就是 \(R\) 本身。
\end{proof}
